% Generated draft: requires mathematical and source review.
\documentclass[11pt]{article}
\usepackage[margin=25mm]{geometry}
\usepackage{fontspec}
\setmainfont{Latin Modern Roman}
\usepackage{amsmath,amssymb,mathtools,unicode-math}
\setmathfont{Latin Modern Math}
\usepackage{enumitem,needspace,xurl,hyperref}
\setlength{\parindent}{0pt}
\setlength{\parskip}{5pt}
\setlength{\emergencystretch}{3em}
\newcommand{\N}{\mathbb{N}}
\newcommand{\Z}{\mathbb{Z}}
\newcommand{\Q}{\mathbb{Q}}
\newcommand{\R}{\mathbb{R}}
\newcommand{\C}{\mathbb{C}}
\newcommand{\F}{\mathbb{F}}
\newcommand{\A}{\mathbb{A}}
\newcommand{\PP}{\mathbb{P}}
\newcommand{\E}{\mathbb{E}}
\newcommand{\eps}{\varepsilon}
\newcommand{\dd}{\,\mathrm{d}}
\begin{document}
\section*{UnsolvedMath ID: 30004057}
\textbf{Source problem code:} catalogue entry
\subsection{Definitions and mathematical statement}
Let \(p\) denote a prime. A unitary Shimura variety means a Shimura variety attached to a Shimura datum whose group is of unitary type. To discuss its modular forms in characteristic \(p\), one must specify an integral model, a characteristic-\(p\) fiber \(S\), and the relevant automorphic bundles. For a specified automorphic bundle \(\mathcal E\) on \(S\), its global sections \(H^0(S,\mathcal E)\) are the corresponding automorphic forms; any conditions at a compactification boundary must also be specified.

The classical structures named in the source are the Hasse invariant, the theta derivation, and theta cycles. In the classical modular-curve setting, the Hasse invariant is a modular form of weight \(p-1\) whose nonvanishing detects the ordinary locus. The classical theta operator is characterized on Fourier expansions by
\[
\theta\!\left(\sum_{n\geq 0}a_nq^n\right)
=q\frac{\dd}{\dd q}\left(\sum_{n\geq 0}a_nq^n\right)
=\sum_{n\geq 0}n a_nq^n.
\]
Theta cycles describe the behavior of successive theta iterates, including their weights or filtrations. This description does not assert periodicity or a particular cycle length in the unitary setting.

\textbf{Requested research program.} Develop, for unitary Shimura varieties, characteristic-\(p\) analogues of the Hasse invariant, the theta derivation, and theta cycles, including the relationships among these structures.

\textbf{Unresolved scope.} The supplied record does not specify the Shimura datum, signature, level, prime hypotheses, integral model, automorphic weights, boundary conditions, or precise properties required of the analogues. It also does not specify whether a Hasse invariant should detect the ordinary locus, the \(\mu\)-ordinary locus, or another stratum, or whether theta operators should be defined globally or on an open locus. Consequently, the source specifies a research program rather than a uniquely quantified existence theorem. No assertion of existence for every unitary Shimura variety or every prime, and no particular weight shift or cycle law, is imposed by this draft.
\subsection{Short English statement}
Develop characteristic-\(p\) versions of Hasse invariants, theta operators, and theta cycles for unitary Shimura varieties. The source leaves the geometric setting and required identities unspecified, so a precise theorem requires further hypotheses. Its supplied \texttt{partially\_solved} label is dated curation, not an independently verified conclusion.
\subsection{Sources}
\begin{itemize}
\item Local registry: \texttt{data/unsolvedmath/problems.json}, permanent UnsolvedMath ID \(30004057\), title ``Theta Operators for Unitary Shimura Varieties,'' exact source problem code \texttt{catalogue entry}; collection: Oberwolfach Reports Open Problems. The supplied record gives proposed year \(2019\) and does not identify an individual proposer.
\item Supplied report citation: \emph{Arithmetic of Shimura Varieties (2020)}, Oberwolfach Reports, DOI \texttt{10.4171/owr/2019/2}, pp.\ 65--131. \url{https://doi.org/10.4171/owr/2019/2}. The supplied background describes these reports as records of workshops and their open-problem discussions. No specific problem page or speaker attribution is supplied.
\item Supplied literature reference: de Shalit--Goren, theta operators. \url{https://doi.org/10.2140/ant.2019.13.1829}.
\item Supplied literature reference: Eischen--Mantovan, entire theta operators. \url{https://arxiv.org/abs/2002.09450}.
\item Supplied literature reference: La Porta, generalized theta operators. \url{https://arxiv.org/abs/2209.03717}.
\item Literature assessment dated exactly \texttt{2026-08-22}: the supplied curation labels the problem partially solved and reports substantial realizations, with theta operators constructed for unitary groups of arbitrary signature in the inert case and more broadly at unramified primes, while a complete theory for all unitary Shimura settings is not established. This is supplied dated curation derived from the OWR record, not an independently verified literature assessment; its time-sensitive status and cited claims require independent verification.
\item The registry labels the statement \texttt{corrected\_verified} and reports a manual semantic, typographic, and LaTeX review confirming preservation of the mathematical target. This is supplied verification metadata, not mathematical review of this draft or independent source verification.
\end{itemize}
\end{document}
